% Part: proof-theory
% Chapter: natural-deduction
% Section: grafting

\documentclass[../../../include/open-logic-section]{subfiles}

\begin{document}

\olfileid[pt]{pt}{nat}{gra}

\olsection{Enxerto de \usetoken{P}{proof}}

!!^{proof}s em \Log{N1c} e~\Log{N1i} são !!{proof}s de !!{formula}s~$!A$ a partir de
hipóteses~$!B$. Suponha que $!B$ tenha, ela própria, !!a{proof}. Como
combinamos a !!{proof} de~$!B$ com a !!{proof} de $!A$ a partir de~$!B$ para
obter !!a{proof} de~$!A$ que não use a hipótese~$!B$?

Uma opção é remover a hipótese~$!B$ da !!{proof} de~$!A$ aplicando
\Intro{\lif}. Podemos então combinar a !!{proof} resultante de
$!B \lif !A$ com a !!{proof} de~$!B$ para obter
\[
\AxiomC{$\Discharge{!B}{x}$}
\DeduceC{$!A$}
\DischargeRule{\Intro{\lif}}{x}
\UnaryInfC{$!B \lif !C$}
\AxiomC{}
\DeduceC{$!B$}
\RightLabel{\Elim{\lif}}
\BinaryInfC{$!A$}
\DisplayProof
\]
Isso é direto, mas um tanto deselegante: por que introduzir $\lif$ apenas
para eliminá-lo imediatamente? Deveria ser possível evitar esse tipo
de ``desvio'' por meio de $!B \lif !A$. (O fato de que esses desvios
sempre podem ser evitados é justamente o conteúdo do teorema da normalização.)

Em \Log{N1c} e~\Log{N1i}, podemos de fato evitá-lo com bastante facilidade neste caso:
basta substituir as hipóteses~$!B^x$ pela !!{proof} inteira
de~$!B$. Esse ``enxerto'' de !!{proof}s nas hipóteses de outras provas
é bastante útil, por isso o registramos em uma definição.

\begin{defn}
Se $\delta_1$ é uma \Log{N1c}-!!{proof} (ou \Log{N1i}-!!{proof})
de~$!A$ a partir da hipótese aberta~$!B^x$, e $\delta_2$ é uma
\Log{N1c}-!!{proof} (\Log{N1i}-!!{proof}) de~$!B$, então
$\Subst{\delta_1}{\delta_2}{!B^x}$ é o resultado de substituir toda
ocorrência não descarregada de~$!B^x$ em~$\delta_1$ por~$\delta_2$.
\end{defn}

Desde que $\delta_1$ e $\delta_2$ não tenham !!{formula}s de hipótese
diferentes marcadas com o mesmo rótulo e nenhuma hipótese aberta
de~$\delta_2$ contenha uma variável própria de~$\delta_1$, a
$\Subst{\delta_1}{\delta_2}{!B^x}$ resultante é uma !!{proof} correta,
e suas hipóteses abertas são as de $\delta_1$ juntamente com as
de~$\delta_2$. Ao enxertar !!{proof}s, essas condições normalmente
serão satisfeitas. A primeira é satisfeita sempre que $\delta_1$ e
$\delta_2$ são subprovas de uma !!{proof} maior. Se a segunda não for
satisfeita, podemos renomear as variáveis próprias de~$\delta_1$ para
que ela seja satisfeita.

Uma definição correspondente se aplica a \Log{N2c} e~\Log{N2i}.

\begin{defn}
Se $\delta_1$ é uma \Log{N2c}-!!{proof} (ou \Log{N2i}-!!{proof})
de~$x:!B, \Gamma_1 \Sequent !A$, e $\delta_2$ é uma
\Log{N1c}-!!{proof} (\Log{N1i}-!!{proof}) de~$\Gamma_2 \Sequent !B$,
então $\Subst{\delta_1}{\delta_2}{x:!B}$ é o resultado de substituir
todo axioma $x:!B \Sequent !B$ em~$\delta_1$ por~$\delta_2$ e acrescentar
$\Gamma_2$ ao contexto de todo sequente abaixo de tais axiomas.
\end{defn}

\begin{prop}
Se $\delta_1 \Proves x:!B, \Gamma_1 \Sequent
!A$ e $\delta_2 \Proves \Gamma_2 \Sequent !B$, então
$\Subst{\delta_1}{\delta_2}{x:!B} \Proves \Gamma_1, \Gamma_2 \Sequent
!A$, desde que nenhuma variável própria de~$\delta_1$ ocorra em~$\Gamma_2$.
\end{prop}


\begin{proof}
  Por indução sobre $\pheight{\delta}$.
\end{proof}

\end{document}
